\section{Results}

\section{Experimental Activity and Analysis of Experimental Results}

\subsection{Preprocessing and Edge Detection comparative analysis}

To empirically validate the optimal preprocessing and feature extraction strategy for the lane detection pipeline, a comparative analysis was conducted across eight distinct bowling video clips. These clips were specifically selected to encompass a variety of challenging real-world conditions, including varying camera exposures, surface glare, and heterogeneous lighting. The study evaluated the combinatorial performance of five single-channel transformations (Default Grayscale, Lightness, Hue, $R+G-B$, and PCA) paired with three standard edge detection algorithms (Sobel, Canny, and Laplacian). A trial was deemed successful strictly if the resulting gradient map enabled the subsequent Hough Transform and heuristic filters to correctly segment the entire playing surface by accurately isolating all four geometric boundaries (the foul line, the pin deck, and both lateral gutters). The quantitative success rates for each combination are detailed in Table \ref{tab:ablation_results}.

An analysis of the aggregated results reveals clear performance disparities not only among the color transformations but, crucially, among the edge detection algorithms. The Sobel operator emerged as the unequivocally strongest edge detection method for correctly identifying lane borders, outperforming Canny and Laplacian across almost every grayscale transformation. In contrast, while the multi-stage Canny detector and the second-order Laplacian filter are exceptionally strong edge detectors, this extreme sensitivity proved detrimental in a visually cluttered bowling alley. Qualitative analysis of the failed clips revealed that Canny and Laplacian frequently detected crisp, high-contrast edges entirely outside the target playing surface, such as architectural lines on the background wall, adjacent lane gutters, or the bowler's silhouette. These highly defined but irrelevant structures generated overwhelmingly strong candidate lines in the Hough space. Consequently, these false positives bypassed the spatial and geometric heuristics of the pipeline, leading to incorrect boundary selection and severely deformed segmentation quadrilaterals (see Figure \ref{fig:over_sensitive_edges}). Sobel, by providing a more subdued and directionally targeted gradient map, allowed the geometric filters to consistently lock onto the true lane boundaries.
 
When observing the single-channel transformations paired with this dominant Sobel filter, a divergence between general cross-filter robustness and peak specific performance became apparent. In terms of general stability across all detectors, the domain-specific $R+G-B$ transformation proved highly robust. It successfully identified boundaries in 5, 4, and 3 clips when paired with the Sobel, Canny, and Laplacian filters, respectively. This consistency validates the underlying hypothesis inherited from autonomous driving literature: mathematically amplifying the red and green channels while suppressing the blue channel effectively isolates the warm, wooden tones of the bowling lane from the cooler, heavily shadowed regions of the gutters and background.

Despite the baseline stability of the $R+G-B$ transformation, the highest success rate was achieved by the standard Default Grayscale conversion paired with the Sobel filter, which successfully localized boundaries in 6 out of the 8 clips. The perceptual luminance weighting of the default grayscale conversion provided the cleanest directional intensity gradients at the lane edges when acted upon by the Sobel operator. Consequently, optimizing for the highest peak accuracy across diverse video inputs, this specific combination was selected as the final implementation for the pipeline.

Finally, the analysis provided valuable insights into the theoretical limitations of chromatic edge detection. It was initially hypothesized that the Hue channel from the HSV color space might yield interesting results by isolating pure color transitions and entirely negating the impact of absolute brightness and shadows. In practice, however, the Hue channel demonstrated the poorest performance of the study, completely failing to localize boundaries across almost all filters. Qualitative analysis of the output (see Figure \ref{fig:hue_failure}) revealed the cause of this failure: the Hue representation effectively erased the critical geometric boundaries (particularly the bottom foul line) from the image entirely. This occurs because the physical demarcations of a bowling lane (such as a dark foul line or shaded gutters contrasting against the light wooden surface) are primarily defined by sharp shifts in luminance and intensity, rather than distinct shifts in core color. By stripping away the intensity information, the Hue channel eliminates the necessary gradient contrast, causing the borders to blend indistinguishably into the background. This visual evidence conclusively confirms that robust boundary extraction in a bowling setting relies fundamentally on structural intensity rather than pure chromaticity.

% --- TABLE OF RESULTS ---
% Note: Make sure to include \usepackage{booktabs} in your document preamble.
\begin{table}[htbp]
\centering
\caption{Quantitative success rates of boundary detection for each preprocessing and edge detection combination across 8 video clips.}
\label{tab:ablation_results}
% This line automatically scales the table to the width of the text:
\resizebox{\textwidth}{!}{%
\begin{tabular}{llccccccccc}
\toprule
\textbf{Grayscale Method} & \textbf{Edge Method} & \textbf{clip1} & \textbf{clip2} & \textbf{clip3} & \textbf{clip4} & \textbf{clip5} & \textbf{clip6} & \textbf{clip7} & \textbf{clip8} & \textbf{Sum} \\
\midrule
default & sobel & 1 & 1 & 1 & 0 & 1 & 1 & 1 & 0 & 6 \\
default & canny & 0 & 0 & 1 & 0 & 1 & 0 & 0 & 0 & 2 \\
default & laplacian & 0 & 0 & 1 & 0 & 1 & 0 & 0 & 0 & 2 \\
\midrule
lightness & sobel & 1 & 0 & 1 & 0 & 1 & 1 & 1 & 0 & 5 \\
lightness & canny & 1 & 0 & 1 & 0 & 0 & 1 & 0 & 0 & 3 \\
lightness & laplacian & 0 & 0 & 1 & 0 & 0 & 1 & 0 & 0 & 2 \\
\midrule
hue & sobel & 0 & 1 & 0 & 0 & 0 & 0 & 0 & 0 & 1 \\
hue & canny & 0 & 0 & 0 & 0 & 0 & 0 & 0 & 0 & 0 \\
hue & laplacian & 0 & 1 & 1 & 0 & 0 & 0 & 0 & 0 & 2 \\
\midrule
r\_g\_minus\_b & sobel & 1 & 1 & 0 & 0 & 1 & 1 & 1 & 0 & 5 \\
r\_g\_minus\_b & canny & 1 & 1 & 1 & 0 & 1 & 0 & 0 & 0 & 4 \\
r\_g\_minus\_b & laplacian & 1 & 1 & 1 & 0 & 0 & 0 & 0 & 0 & 3 \\
\midrule
pca & sobel & 1 & 0 & 1 & 0 & 1 & 1 & 1 & 0 & 5 \\
pca & canny & 1 & 0 & 1 & 0 & 0 & 1 & 0 & 0 & 3 \\
pca & laplacian & 0 & 0 & 1 & 0 & 1 & 0 & 1 & 0 & 3 \\
\bottomrule
\end{tabular}%
} % Closes the resizebox
\end{table}

% --- FIGURES ---
\begin{figure}[htbp]
    \centering
    % Replace the placeholder file names with your actual image paths
    % \includegraphics[width=0.48\textwidth]{images/default_canny_quad.png}
    % \hfill
    % \includegraphics[width=0.48\textwidth]{images/pca_laplacian_quad.png}
    \caption{Examples of heuristic failure due to over-sensitive edge detection. Left: Canny detects background wall structures. Right: Laplacian detects the bowler's leg and adjacent lane. Both result in an incorrect homography quadrilateral.}
    \label{fig:over_sensitive_edges}
\end{figure}

\begin{figure}[htbp]
    \centering
    % \includegraphics[width=0.6\textwidth]{images/hue_laplacian_bottom_segments.png}
    \caption{Output of the Hue channel on the lower lane segment. By discarding luminance data, the critical intensity contrast of the foul line is completely erased, resulting in a failure to detect the bottom boundary.}
    \label{fig:hue_failure}
\end{figure}




\subsection{Top Boundary Localization via Object-Centric Regression}

As established in the methodology, the top boundary of the lane (the pin deck interface) presents a unique geometric challenge: the physical edge is heavily occluded by the bowling pins themselves, making standard gradient-based edge detection ineffective. To overcome this occlusion, an object-centric approach was evaluated, employing multi-scale template matching to detect the pins and utilizing the bottom-center coordinates of their bounding boxes as spatial candidate points for the boundary. The specific selection of the bottom-center coordinate, rather than the geometric centroid (midpoint) of the bounding box, is a critical design choice. Empirical analysis of the bounding box outputs revealed that using the centroid artificially elevated the regression line, mapping it across the midsections of the pins rather than the lane surface. Because the physical objective is to map the pin deck interface, the bottom-center coordinate provides a geometrically accurate anchor directly to the true floor plane.

While multi-scale template matching successfully localized the target pins, the raw output was inherently noisy. As anticipated in a visually complex bowling environment, the algorithm frequently generated false positive detections, predominantly from pin reflections on the highly polished synthetic oil and visually similar architectural features in the background. Interestingly, the template matching also frequently detected pins actively standing in adjacent parallel lanes. Rather than treating these adjacent detections as noise, the pipeline exploits them to increase the stability of the boundary estimation. Because bowling alleys are constructed with coplanar, horizontally aligned pin decks, pins in adjacent lanes share the exact same geometric baseline as the target lane. Consequently, these adjacent detections serve as valid, extended data points that drastically improve the stability of the linear estimation.

To robustly extract the true boundary from this spatial data while filtering out the actual noise, a Random Sample Consensus (RANSAC) regressor was evaluated. RANSAC is uniquely suited for this task due to its outlier rejection capabilities. By iteratively estimating a first-degree polynomial model from random subsets of the spatial coordinates, the algorithm successfully identified the dominant linear structure (the continuous base of the pins extending across the visible lanes) while ignoring coordinate data from reflections and background architecture that did not conform to this geometry.

Qualitative analysis of the regression output (see Figure \ref{fig:ransac_pin_deck}) demonstrates the high efficacy of this approach across varying camera angles and lane designs. The RANSAC algorithm successfully segmented the raw candidate points into a strict inlier consensus set (visualized as blue points) and a rejected outlier set (visualized as red points). The inlier set actively incorporates valid pins from both the target and adjacent lanes, exploiting the extended horizontal baseline to anchor the regression more firmly. The red spatial outliers correctly encapsulate the false detections caused by reflections and background noise. By fitting the final line strictly to these fortified blue inliers, the pipeline generated a highly accurate and stable top boundary line (visualized in green). This confirms that decoupling the structural detection into an object-centric model effectively bypasses severe visual occlusions at the pin deck. 

% --- FIGURE FOR RANSAC ---
\begin{figure}[htbp]
    \centering
    % Ensure you have \usepackage{graphicx} in your preamble
    % Replace the placeholder below with the actual path to your image
    % \includegraphics[width=0.9\textwidth]{images/top_boundary_midpoints.jpg}
    \caption{Performance of the RANSAC regressor on multi-scale template matching output. The bottom-center coordinates of all detected bounding boxes are plotted. The algorithm successfully classifies true pins (including those actively exploited on adjacent, aligned lanes) as inliers (blue dots). It effectively rejects false positives from reflections and background structures as outliers (red dots). The final fitted polynomial (green line) accurately maps the occluded pin deck boundary using the extended geometric baseline.}
    \label{fig:ransac_pin_deck}
\end{figure}